\documentclass[UTF8,12pt]{ctexart}
\usepackage[a4paper,margin=24mm]{geometry}
\usepackage{amsmath,amssymb,bm,graphicx,adjustbox}
\newcommand{\fitmath}[1]{\begin{adjustbox}{max width=\linewidth}$\displaystyle #1$\end{adjustbox}}
\setlength{\parindent}{0pt}
\setlength{\parskip}{.7em}
\sloppy
\allowdisplaybreaks
\begin{document}
\section*{1990 年考研数学三 · 参考答案}
\subsection*{1990年真题参考答案}

\subsection*{（试卷IV）}

\subsection*{一、填空题}

（1）\(\displaystyle 2\)。 （2）\(\displaystyle a+b\)。 （3）\(\displaystyle \dfrac{\displaystyle 9}{\displaystyle 2}\)。 （4）\(\displaystyle a_1+a_2+a_3+a_4=0\)。 （5）\(\displaystyle \dfrac{\displaystyle 2}{\displaystyle 3}\)。


\subsection*{二、选择题}

（1）B。 （2）D。 （3）C。 （4）A。 （5）C。


\subsection*{三、计算题}

（1）\(\displaystyle \ln(1+e)-\dfrac{\displaystyle e}{\displaystyle 1+e}\)。 （2）\(\displaystyle \dfrac{\displaystyle 5}{\displaystyle 144}\)。 （3）\(\displaystyle [2,\allowbreak 4]\)。 （4）\(\displaystyle y=e^{-\sin x}(x\ln x-x+C)\)，其中 \(\displaystyle C\) 为任意常数。


四、（1）当 \(\displaystyle x_1=0.75,\allowbreak x_2=1.25\) 时，利润取得最大值。（2）当仅有1.5万元广告费用时，最优广告策略是将1.5万元费用全部用于报纸的广告费。


五、证明略（在 \(\displaystyle [0,\allowbreak a]\) 和 \(\displaystyle [b,\allowbreak a+b]\) 上分别应用拉格朗日中值定理）。


六、（1）\(\displaystyle a=1,\allowbreak b=3\)。 （2）基础解系为 \(\displaystyle \boldsymbol\xi_1=(1,\allowbreak -2,\allowbreak 1,\allowbreak 0,\allowbreak 0)^{\mathrm T},\allowbreak \ \boldsymbol\xi_2=(1,\allowbreak -2,\allowbreak 0,\allowbreak 1,\allowbreak 0)^{\mathrm T},\allowbreak \ \boldsymbol\xi_3=(5,\allowbreak -6,\allowbreak 0,\allowbreak 0,\allowbreak 1)^{\mathrm T}\)。 （3）方程的全部解为 \(\displaystyle \boldsymbol x=(-2,\allowbreak 3,\allowbreak 0,\allowbreak 0,\allowbreak 0)^{\mathrm T}+k_1\boldsymbol\xi_1+k_2\boldsymbol\xi_2+k_3\boldsymbol\xi_3\)，其中 \(\displaystyle k_1,\allowbreak k_2,\allowbreak k_3\) 为任意常数。


七、证明略。\(\displaystyle (E-A)^{-1}=E+A+\cdots+A^{k-1}\)。


八、证明略（反证法）。


九、\(\displaystyle P(A_1)=\dfrac{\displaystyle 7}{\displaystyle 15};\ P(A_2)=\dfrac{\displaystyle 14}{\displaystyle 15};\ P(A_3)=\dfrac{\displaystyle 7}{\displaystyle 30}\)。


十、（1）\(\displaystyle X\) 和 \(\displaystyle Y\) 独立。（2）\(\displaystyle e^{-0.1}\)。


十一、\(\displaystyle 0.682\)。


\subsection*{（试卷V）}

\subsection*{一、填空题}

（1）【同试卷IV 第一、（1）题】 （2）【同试卷IV 第一、（2）题】 （3）【同试卷IV 第一、（3）题】 （4）【同试卷IV 第一、（4）题】 （5）\(\displaystyle N(0,\allowbreak 5)\)。


\subsection*{二、选择题}

（1）【同试卷IV 第二、（1）题】 （2）【同试卷IV 第二、（2）题】 （3）【同试卷IV 第二、（3）题】 （4）A。 （5）B。


\subsection*{三、计算题}

（1）\(\displaystyle \dfrac{\displaystyle 1}{\displaystyle 2}\)。 （2）\(\displaystyle -\dfrac{\displaystyle 1}{\displaystyle 8}x\csc^2\dfrac{\displaystyle x}{\displaystyle 2}-\dfrac{\displaystyle 1}{\displaystyle 4}\cot\dfrac{\displaystyle x}{\displaystyle 2}+C\)，其中 \(\displaystyle C\) 为任意常数。


\subsection*{（试卷V）参考答案续}

三、（3）\(\displaystyle \dfrac{\displaystyle \partial z}{\displaystyle \partial y}=\dfrac{\displaystyle y\varphi(z/y)-z\varphi\prime(z/y)}{\displaystyle 2yz-y\varphi\prime(z/y)}\)。 （4）【同试卷IV 第三、（2）题】


四、【同试卷IV 第四题】


五、证明略（可考虑函数 \(\displaystyle f(x)=1+x\ln(x+\sqrt{1+x^2})-\sqrt{1+x^2}\)，计算 \(\displaystyle f\prime(x)\)，证明 \(\displaystyle f(x)\) 的最小值为 \(\displaystyle f(0)=0\)）。


六、\(\displaystyle \lambda^{10}-10^{10}\)。


七、证明略（考虑 \(\displaystyle (A\boldsymbol x)^{\mathrm T}(A\boldsymbol x)\)，利用 \(\displaystyle A^{\mathrm T}A=E\) 可得 \(\displaystyle (A\boldsymbol x)^{\mathrm T}(A\boldsymbol x)=\boldsymbol x^{\mathrm T}\boldsymbol x\)）。


八、【同试卷IV 第六题】 九、【同试卷IV 第九题】


十、


\[\fitmath{\displaystyle \begin{array}{c|ccc}Y\backslash X&0&1&2\\\hline0&0.16&0.08&0.01\\1&0.32&0.16&0.02\\2&0.16&0.08&0.01\end{array}}\]


十一、【同试卷IV 第十一题】

\end{document}
